\section{Utility invariance}
Write $G_e,G_c,G_o$ for the exact, cardinal and ordinal automorphism groups.
\begin{proposition}[U2: canonical representations]
These are subgroups of the same typed-bijection group, with
$G_e\le G_c\le G_o$. Let $c_i=(u_i-\min u_i)/(\max u_i-\min u_i)$ for
nonconstant $u_i$, and $c_i=0$ otherwise. Let $r_i(s)$ be the zero-based rank
of $u_i(s)$ among its distinct levels. Then
\[
G_c=\Aut_e(c),\qquad G_o=\Aut_e(r).
\]
\end{proposition}
\begin{proof}
Affine relations compose with positive slopes and invert; finite-range order
bijections also compose and invert. A cardinal map transports extrema by the
same affine relation and hence preserves $c$. Conversely equality of normalized
tables reconstructs a positive affine relation, including the constant case.
An increasing finite-range bijection carries the $k$th distinct level to the
$k$th distinct level, giving the rank identity and its converse.
\end{proof}

The inclusions can be strict even for player transitivity. For binary actions
and $k=s_0+s_1+s_2$, take $u_i=(i+1)k^2$. Every bare player permutation is
cardinal, but distinct numerical ranges prevent exact player movement.
Alternatively, $u=(k,k^2,k)$ has full ordinal player image, while incompatible
normalized level spacings prevent cardinal maps moving player 1.

\begin{proposition}[U3: focal-fixed invariance]
An ordinal automorphism fixing player $i$ preserves $u_i$ exactly. Consequently
$\ker(q|G_e)=\ker(q|G_c)=\ker(q|G_o)$.
If $G_e$ is player transitive, then $G_e=G_c=G_o$.
\end{proposition}
\begin{proof}
An increasing permutation of one finite ordered payoff range is the identity.
This proves the first assertion and the kernel equality. Exact transitivity
makes all players' numerical ranges identical. Any ordinal transport between
these identical ordered lists must then preserve every level exactly.
\end{proof}
The normalizations represent preferences; they do not calibrate a participant's
cross-role utility or justify expected rank as welfare.
